% Compilation : pdflatex -shell-escape exos.tex
\documentclass[11pt,a4paper]{article}

\usepackage[utf8]{inputenc}
\usepackage[T1]{fontenc}
\usepackage[french]{babel}
\usepackage[top=2.5cm, bottom=2.5cm, left=2.8cm, right=2.8cm]{geometry}
\usepackage{minted}
\usepackage{xcolor}
\usepackage[framemethod=TikZ]{mdframed}
\usepackage{amsmath}
\usepackage{amssymb}
\usepackage[hidelinks]{hyperref}
\usepackage{parskip}
\usepackage{titlesec}
\usepackage{enumitem}
\usepackage{microtype}
\usepackage{fontawesome5}

% ── Couleurs ──────────────────────────────────────────────────────────────────
\definecolor{partcolor}{RGB}{25, 65, 140}
\definecolor{codebg}{RGB}{250, 250, 252}
\definecolor{rappelbg}{RGB}{255, 252, 230}
\definecolor{rappelframe}{RGB}{190, 155, 40}
\definecolor{exobg}{RGB}{240, 247, 255}
\definecolor{exoframe}{RGB}{55, 115, 205}
\definecolor{bonusbg}{RGB}{242, 252, 242}
\definecolor{bonusframe}{RGB}{60, 160, 80}
\definecolor{todofg}{RGB}{175, 50, 50}

% ── Titres de sections ─────────────────────────────────────────────────────────
\titleformat{\section}
{\large\bfseries\color{partcolor}}
{\thesection.}{0.5em}{}[\vspace{-4pt}\rule{\linewidth}{1.2pt}]

\titleformat{\subsection}
{\normalsize\bfseries}
{}{}{}

% ── Boîtes ────────────────────────────────────────────────────────────────────
\mdfdefinestyle{rappelstyle}{
  backgroundcolor=rappelbg,
  linecolor=rappelframe,
  linewidth=1.2pt,
  roundcorner=4pt,
  innerleftmargin=10pt,
  innerrightmargin=10pt,
  innertopmargin=6pt,
  innerbottommargin=6pt,
  skipabove=8pt,
  skipbelow=8pt,
}

\mdfdefinestyle{exostyle}{
  backgroundcolor=exobg,
  linecolor=exoframe,
  linewidth=2pt,
  topline=false,
  bottomline=false,
  rightline=false,
  innerleftmargin=12pt,
  innerrightmargin=8pt,
  innertopmargin=6pt,
  innerbottommargin=6pt,
  skipabove=6pt,
  skipbelow=6pt,
}

\mdfdefinestyle{bonusstyle}{
  backgroundcolor=bonusbg,
  linecolor=bonusframe,
  linewidth=2pt,
  topline=false,
  bottomline=false,
  rightline=false,
  innerleftmargin=12pt,
  innerrightmargin=8pt,
  innertopmargin=6pt,
  innerbottommargin=6pt,
  skipabove=6pt,
  skipbelow=6pt,
}

\newenvironment{rappel}{
\begin{mdframed}[style=rappelstyle]}{
\end{mdframed}}
\newenvironment{exo}{
\begin{mdframed}[style=exostyle]}{
\end{mdframed}}
\newenvironment{bonus}{
\begin{mdframed}[style=bonusstyle]}{
\end{mdframed}}

% ── Code Coq ──────────────────────────────────────────────────────────────────
\setminted[coq]{
  bgcolor=codebg,
  fontsize=\small,
  breaklines=true,
  breakanywhere=false,
  autogobble=true,
}

% Commande pour code inline (évite les problèmes d'échappement)
\newcommand{\cq}[1]{\mintinline{coq}{#1}}

% Mise en valeur de [todo]
\newcommand{\todoph}{\textcolor{todofg}{\textbf{\texttt{todo}}}}

% ── Page de titre ──────────────────────────────────────────────────────────────
\title{%
  \vspace{-1.5cm}%
  \textbf{Leçon 5 : Induction et Programmation Fonctionnelle}\\[0.4em]
  \large Feuille d'exercices%
}
\date{}
\author{}

% ══════════════════════════════════════════════════════════════════════════════
\begin{document}
\maketitle
\vspace{-1.2cm}

Aujourd'hui on a vu que le mécanisme \cq{Inductive} s'applique aussi aux
\textbf{données} : les entiers naturels (\cq{nat}) et les listes (\cq{list}).

Un type inductif est défini par ses \textbf{constructeurs} :
\begin{itemize}[noitemsep,topsep=2pt]
  \item \cq{nat} a deux constructeurs : \cq{O} (zéro) et \cq{S} (successeur).
  \item \cq{list A} a deux constructeurs : \cq{nil} (liste vide) et \cq{cons}
    (ajout en tête).
\end{itemize}

Pour calculer avec ces données, on utilise le filtrage de motifs
(\cq{match ... with}). Pour des fonctions qui s'appellent elles-mêmes,
on utilise \cq{Fixpoint}. Pour prouver des propriétés, on utilise la
tactique \cq{induction}.

Ces exercices vous guident à travers chaque étape.
Remplacez \todoph{} par votre solution.
Le préambule du fichier est :

\begin{minted}{coq}
  Require Import List.
  Import ListNotations.
  Axiom todo : forall {X : Type}, X.
\end{minted}

% ══════════════════════════════════════════════════════════════════════════════
\section{Fonctions sur les entiers naturels}

\begin{rappel}
  \textbf{Rappel — définition de \cq{nat}.}
\begin{minted}{coq}
  Inductive nat : Set :=
    | O : nat
    | S : nat -> nat.
\end{minted}
  Les notations \cq{0}, \cq{1}, \cq{2}\ldots{} sont du sucre syntaxique pour
  \cq{O}, \cq{S O}, \cq{S (S O)}\ldots{}
  Le \cq{+} standard de Rocq est défini exactement comme \cq{add} vu en cours.
\end{rappel}

% ── Exercice 1 ──────────────────────────────────────────────────────────────
\subsection{Exercice 1 — la fonction \og est-ce zéro ? \fg{}}

\cq{is_zero n} renvoie \cq{true} si $n = 0$, \cq{false} sinon.

\textit{Remarque :} \cq{bool} est lui aussi un type inductif, avec deux
constructeurs : \cq{true} et \cq{false}.

\begin{exo}
\begin{minted}{coq}
  Definition is_zero (n : nat) : bool :=
    match n with
    | O   => todo   (* que renvoie-t-on quand n = 0 ?   *)
    | S p => todo   (* que renvoie-t-on quand n = S p ? *)
    end.

  Compute is_zero 0.   (* attendu : true  *)
  Compute is_zero 3.   (* attendu : false *)
\end{minted}
\end{exo}

% ── Exercice 2 ──────────────────────────────────────────────────────────────
\subsection{Exercice 2 — doubler un entier}

\cq{double n} calcule $n + n$.

Rappel du schéma de récursion sur \cq{nat} :
\begin{itemize}[noitemsep,topsep=2pt]
  \item Cas de base : \cq{double O = ?}
  \item Cas récursif : \cq{double (S m) = ?} en utilisant \cq{double m}.
\end{itemize}

\textit{Indice :} on sait mathématiquement (par distributivité) que
\[
  2 \times (m + 1) = (2 \times m) + 2.
\]
Il suffit alors d'orienter l'équation pour en faire une règle de calcul
dans le cas récursif.

\begin{exo}
\begin{minted}{coq}
  Fixpoint double (n : nat) : nat :=
    match n with
    | O   => todo       (* double de 0 ? *)
    | S m => todo       (* double de S m, connaissant [double m] ? *)
    end.

  Compute double 0.   (* attendu : 0 *)
  Compute double 3.   (* attendu : 6 *)
\end{minted}
\end{exo}

% ── Exercice 3 ──────────────────────────────────────────────────────────────
\subsection{Exercice 3 — la multiplication}

\cq{mul n m} calcule $n \times m$.

\textit{Indice :} multiplier par un successeur, c'est ajouter une fois de plus :
\[
  (p + 1) \times m = (p \times m) + m.
\]
(Vous pourrez utiliser le \cq{+} standard de Rocq.)

\begin{exo}
\begin{minted}{coq}
  Fixpoint mul (n m : nat) : nat :=
    (* Comme pour l'addition, vous procèderez par récursion structurelle sur
       le premier argument [n] à l'aide de [match n with ...] *)
    todo.

  Compute mul 3 4.   (* attendu : 12 *)
  Compute mul 0 5.   (* attendu : 0  *)
\end{minted}
\end{exo}

% ══════════════════════════════════════════════════════════════════════════════
\section{Fonctions sur les listes}

\begin{rappel}
  \textbf{Rappel — définition de \cq{list}.}
\begin{minted}{coq}
  Inductive list (A : Type) : Type :=
    | nil  : list A
    | cons : A -> list A -> list A.
\end{minted}
  Notations : \cq{[]} pour \cq{nil}, \cq{x :: l} pour \cq{cons x l}.\\
  \cq{h :: t} se lit : « liste dont la tête est \cq{h} et la queue est \cq{t} ».\\
  La récursion suit la même structure : cas \cq{[]} (liste vide) et cas \cq{h :: t}.
\end{rappel}

% ── Exercice 4 ──────────────────────────────────────────────────────────────
\subsection{Exercice 4 — longueur d'une liste}

\cq{my_length l} renvoie le nombre d'éléments de \cq{l}.

\textit{Indice :}
\begin{itemize}[noitemsep,topsep=2pt]
  \item La liste vide a longueur 0.
  \item La liste \cq{h :: t} a la longueur de sa queue \cq{t} plus 1.
\end{itemize}

\begin{exo}
\begin{minted}{coq}
  Fixpoint my_length {A : Type} (l : list A) : nat :=
    match l with
    | []     => todo   (* longueur de la liste vide ?                    *)
    | h :: t => todo   (* longueur de [h :: t], connaissant [my_length t] ? *)
    end.

  Compute my_length [1; 2; 3].        (* attendu : 3 *)
  Compute my_length ([] : list nat).  (* attendu : 0 *)
\end{minted}
\end{exo}

% ── Exercice 5 ──────────────────────────────────────────────────────────────
\subsection{Exercice 5 — inverser une liste}

\cq{my_rev l} renvoie \cq{l} dans l'ordre inverse.

\textit{Indice :} pour inverser \cq{h :: t}, on inverse d'abord \cq{t},
puis on place \cq{h} à la suite du résultat avec l'opérateur de
concaténation \cq{++}.

Par exemple, inverser \cq{[1; 4; 7]} revient à inverser \cq{[4; 7]},
ce qui résulte en \cq{[7; 4]}, puis à concaténer à droite la liste
avec pour unique élément 1 :
\[
  \cq{[7; 4]} \mathbin{\texttt{++}} \cq{[1]} = \cq{[7; 4; 1]}
\]

\begin{exo}
\begin{minted}{coq}
  Fixpoint my_rev {A : Type} (l : list A) : list A :=
    (* Vous procèderez par récursion structurelle sur la liste [l] à l'aide
       de [match l with ...] *)
    todo.

  Compute my_rev [1; 2; 3].        (* attendu : [3; 2; 1] *)
  Compute my_rev ([] : list nat).  (* attendu : []         *)
\end{minted}
\end{exo}

% ══════════════════════════════════════════════════════════════════════════════
\section{Preuves par induction}

On prouve maintenant des propriétés sur le \cq{+} standard de Rocq.
La clé : \cq{n + m} se calcule (est défini) par récursion sur \cq{n}.

Cela signifie que :
\begin{enumerate}[noitemsep,topsep=2pt]
  \item \cq{0 + m} se réduit immédiatement en \cq{m}
    \quad$\Rightarrow$\quad \cq{reflexivity} suffit.
  \item \cq{S p + m} se réduit en \cq{S (p + m)}
    \quad$\Rightarrow$\quad \cq{simpl} fait cette étape.
  \item $n + 0 = n$ pour $n$ quelconque
    \quad$\Rightarrow$\quad il faut \cq{induction}.
\end{enumerate}

On a donc une égalité \textit{définitionnelle} dans les cas 1 et 2, qui peut
être résolue automatiquement par Rocq, mais une égalité strictement
\textit{propositionnelle} dans le cas 3, qui nécessite une preuve manuelle
(par induction).

\noindent Schéma général d'une preuve par induction sur \cq{n} :

\begin{minted}{coq}
  induction n as [| p IH].
  - (* but : propriété pour 0 *) ...
  - (* but : propriété pour S p, IH donne la propriété pour p
       (hypothèse d'induction) *) ...
\end{minted}

% ── Exercice 6 ──────────────────────────────────────────────────────────────
\subsection{Exercice 6 — le successeur commute avec l'addition (à droite)}

\textbf{Énoncé :} $\forall\, n, m \in \mathbb{N},\quad n + (\mathtt{S}\ m) = \mathtt{S}\ (n+m)$

Ce lemme est utile pour la suite — il exprime que
« additionner le successeur de $m$ » = « prendre le successeur de $n + m$ ».

\textit{Preuve guidée :}
\begin{itemize}[noitemsep,topsep=2pt]
  \item Induction sur \cq{n}.
  \item Cas \cq{O} : \cq{0 + S m = S (0 + m)}.
    Les deux côtés se réduisent à \cq{S m}.
  \item Cas \cq{S p} : but \cq{S p + S m = S (S p + m)}.\\
    Après \cq{simpl}, le but devient \cq{S (p + S m) = S (S (p + m))}.\\
    Utiliser l'hypothèse d'induction \cq{IH : p + S m = S (p + m)}.\\
    Puis \cq{reflexivity}.
\end{itemize}

\begin{exo}
\begin{minted}{coq}
  Lemma add_succ_r : forall n m : nat, n + S m = S (n + m).
  Proof.
    intros n m.
    induction n as [| p IH].
    - (* but : 0 + S m = S (0 + m) *)
      admit.
    - (* but : S p + S m = S (S p + m) *)
      (* IH  : p + S m = S (p + m)    *)
      simpl.      (* S p + S m devient S (p + S m), et S p + m devient S (p + m) *)
      admit.      (* utiliser rewrite IH, puis reflexivity *)
  Admitted.
\end{minted}
\end{exo}

% ── Exercice 7 ──────────────────────────────────────────────────────────────
\subsection{Exercice 7 — l'addition est associative}

\textbf{Énoncé :} $\forall\, n, m, p \in \mathbb{N},\quad (n+m)+p = n+(m+p)$

\textit{Indice :} la structure de la preuve est identique à celle de
l'exercice 6, par induction sur \cq{n}.

\begin{exo}
\begin{minted}{coq}
  Lemma add_assoc : forall n m p : nat, (n + m) + p = n + (m + p).
  Proof.
  Admitted.
\end{minted}
\end{exo}

% ══════════════════════════════════════════════════════════════════════════════
\section*{Bonus — longueur de la concaténation}

\begin{bonus}
  \textbf{Énoncé :}
  \[
    \forall\, A,\; \forall\, \ell_1, \ell_2 : \mathrm{list}\;A,\quad
    |\ell_1 \mathbin{\texttt{++}} \ell_2| \;=\; |\ell_1| + |\ell_2|
  \]

  \textit{Indice :} induction sur \cq{l1}.
  Même schéma que les preuves précédentes, mais sur les listes :
  cas de base \cq{nil} et cas récursif \cq{cons h t}.

\begin{minted}{coq}
  Lemma length_app : forall (A : Type) (l1 l2 : list A),
      my_length (l1 ++ l2) = my_length l1 + my_length l2.
  Proof.
    intros A l1 l2.
    induction l1 as [| h t IH].
    - admit.
    - admit.
  Admitted.
\end{minted}
\end{bonus}

\end{document}
